% 24-21(24쪽) — 오른쪽 컬럼 공통 그림: 원 x^2+y^2=1 위의 점 P(t, √(1-t²)) 와 점 A(1, 0) 을 지나는 직선이 y축과 만나는 점 B · P 에서 x축에 내린 수선의 발 H. 스캔 화질이 낮아 TikZ 로 다시 그림.
% 원문에 있는 것만: 눈금 1, -1(두 축) · 라벨 x^2+y^2=1, B, P(t, √(1-t²)), A, H, O, x, y 와 H 의 직각 표시.
\begin{tikzpicture}[x=1.15cm, y=1.15cm, line width=0.5pt]
  \draw[->, >=stealth, line width=0.7pt] (-1.95,0) -- (2.15,0) node[below=1pt] {$x$};
  \draw[->, >=stealth, line width=0.7pt] (0,-1.80) -- (0,2.50) node[right=1pt] {$y$};
  \draw[line width=0.6pt] (0,0) circle (1);
  % 두 점 A(1, 0), P(t, √(1-t²)) 를 지나는 직선 (y축과 만나는 점이 B)
  \draw[line width=0.5pt] (-0.30,2.252) -- (1.85,-1.472);
  % 수선 PH
  \draw[line width=0.5pt] (0.5,0.8660) -- (0.5,0);
  % H 의 직각 표시
  \draw[line width=0.4pt] (0.36,0) -- (0.36,0.14) -- (0.50,0.14);
  \fill (0.5,0.8660) circle (2.2pt);
  \node[anchor=east, font=\small] at (-1.05,1.00) {$x^2+y^2=1$};
  \node[anchor=west, font=\small] at (0.60,1.00) {$\pt{P}(t,\,\sqrt{1-t^2})$};
  \node[anchor=east, font=\small, inner sep=2pt] at (-0.02,1.732) {$\pt{B}$};
  \node[anchor=south west, font=\small, inner sep=1pt] at (1.05,0.05) {$\pt{A}$};
  \node[below left=-2pt and -3pt, font=\small] at (0,0) {$\pt{O}$};
  \node[below=1pt, font=\small] at (0.50,0) {$\pt{H}$};
  \node[below=1pt, font=\small] at (0.90,0) {$1$};
  \node[below=1pt, font=\small] at (-1.00,0) {$-1$};
  \node[anchor=east, font=\small, inner sep=2pt] at (-0.02,1.00) {$1$};
  \node[anchor=east, font=\small, inner sep=2pt] at (-0.02,-1.00) {$-1$};
\end{tikzpicture}
